% Figure 2.1: a question at a glance, from the chat to the answer (details in the next figure)
\begin{figure}[H]
  \centering
  \begin{tikzpicture}[font=\scriptsize,
      lay/.style={rounded corners=3pt, text width=1.5cm, minimum height=1.55cm, align=center,
                  inner sep=2pt},
      flow/.style={-{Stealth[length=1.8mm]}, accent, line width=0.7pt}]
    \newcommand\lgx[1]{\includegraphics[height=0.5cm, width=0.62cm, keepaspectratio]{logos/#1}}
    % \lay{index}{category}{icon or logo}{label}
    \newcommand\lay[4]{\node[lay, cat #2] (l#1) at ({(#1-1)*1.98},0)
      {#3\\[3pt]{\bfseries #4}};}
    \lay{1}{infra}{{\color{accent}\Large\faComments}}{Web chat}
    \lay{2}{sec}{\lgx{keycloak}}{Optional login}
    \lay{3}{sec}{\lgx{qwen}}{Input guardrail}
    \lay{4}{infra}{\lgx{langgraph}}{Orche\-strator}
    \lay{5}{data}{\lgx{opensearch}\,\lgx{mcp}}{Sources}
    \lay{6}{model}{\lgx{openai}}{Generation}
    \lay{7}{sec}{{\color{catsec}\Large\faCheckDouble}}{Output guardrail}
    \lay{8}{infra}{\lgx{chatwoot}}{Answer or operator}
    \foreach \a/\b in {1/2,2/3,3/4,4/5,5/6,6/7,7/8}
      \draw[flow] (l\a) -- (l\b);
  \end{tikzpicture}
  \caption{A question at a glance. Figure~\ref{fig:architecture} shows each layer in detail.}
  \label{fig:overview-strip}
\end{figure}
